\section{Excluding five or more active cells}\label{elim:section}

We show that the scalar bounds of Lemma~\ref{sc:bounds} and the
geometric constraints of Corollary~\ref{geo:constraints} are incompatible.
We first confine the smallest radius to $.066<t<.12$, and then prove
$c>2.6t$, contradicting the second determinant constraint.
Throughout this section, the radii are ordered as
\[
 r_1\ge\cdots\ge r_m=t>0,\qquad
 \sum_{j=1}^m r_j^2=1,\qquad
 L=r_1^2+r_2^2,\qquad c=r_3^2.
\]
Set
\begin{equation}\label{elim:parameters}
 A=\begin{cases}.929,&m=5,\\.884,&m\ge6,\end{cases}
 \qquad d=A-.15,\qquad s=\frac{.415}{d},
 \qquad e(u)=.415u-du^2.
\end{equation}
The probability excesses
\(\delta_j=P_j-Ar_j^2\) satisfy, by \eqref{sc:three-bounds},
\begin{equation}\label{elim:budget}
 \delta_j\ge\max\{0,e(r_j)\},\qquad
 \sum_{j=1}^m\delta_j=1-A.
\end{equation}
All terminating decimals below denote exact rational numbers. The finite
comparisons are recorded with strict lower bounds; the rational arithmetic
and square-root enclosures that verify them are described in
Appendix~\ref{cert:arithmetic}.

\subsection{The third radius and the small-minimum case}

First we obtain a bound that holds for every value of~$t$.
For $j\ge3$, the inequality $r_j\le\sqrt c$ gives
\[
 e(r_j)=r_j^2\left(\frac{.415}{r_j}-d\right)
 \ge r_j^2\left(\frac{.415}{\sqrt c}-d\right).
\]
Summing these charges in \eqref{elim:budget} yields
\[
 (1-L)\left(\frac{.415}{\sqrt c}-d\right)\le1-A.
\]
If $A=.929$, add $.045(1-L)$ to both sides and bound the added
right-hand term by $.045$; if $A=.884$, no adjustment is needed. Thus in both cases
\begin{equation}\label{elim:common-charge}
 (1-L)\left(\frac{.415}{\sqrt c}-.734\right)\le.116.
\end{equation}
Since $r_j\le2/3$, we have $L\le8/9$. If $\sqrt c\le.23$,
the left side of \eqref{elim:common-charge} exceeds its right side by
at least
\[
 \frac19\left(\frac{.415}{.23}-.734\right)-.116
 =\frac{101}{34500}>0.
\]
Consequently
\begin{equation}\label{elim:third-radius}
 \sqrt c>.23.
\end{equation}

Suppose now that $t\le q_0:=.066$. We contradict
\eqref{geo:first-constraint}, which reads
\begin{equation}\label{elim:first-constraint}
 (c-t^2)(L-c-2t^2)\le8(3/\pi)^2t^2.
\end{equation}
If $\sqrt c\ge.44$, then $L\ge2c$, so the left side is at least
\[
 (c-t^2)(c-2t^2)
 \ge(.44^2-q_0^2)(.44^2-2q_0^2).
\]
Both factors here are positive, and
\begin{equation}\label{elim:large-third-small-minimum}
 (.44^2-q_0^2)(.44^2-2q_0^2)
       -8(3/\pi)^2q_0^2>.00321132.
\end{equation}
This contradicts \eqref{elim:first-constraint}.

For the remaining range $.23<\sqrt c<.44$, use the four intervals
in Table~\ref{elim:small-table}. On an interval $[a,b]$, the denominator
$.415/\sqrt c-.734$ is positive, and
\eqref{elim:common-charge} gives
\[
 L-c-2t^2
 \ge 1-\frac{.116}{.415/b-.734}-b^2-2q_0^2
 =:M(b).
\]
Also $c-t^2\ge a^2-q_0^2>0$. Every $M(b)$ is positive, and the
last column of the table contradicts \eqref{elim:first-constraint}.
Thus
\begin{equation}\label{elim:minimum-lower}
 t>.066.
\end{equation}

\begin{table}[ht]
\centering
\begin{tabular}{c|c|c}
$[a,b]$ & \shortstack{Lower bound for\\ $M(b)$} &
 \shortstack{Lower bound for\\ $(a^2-q_0^2)M(b)-8(3/\pi)^2q_0^2$}\\
\hline
$[.23,.30]$ & $.72264$ & $.00330237$\\
$[.30,.35]$ & $.61198$ & $.02063552$\\
$[.35,.41]$ & $.40621$ & $.01621418$\\
$[.41,.44]$ & $.24314$ & $.00803616$
\end{tabular}
\caption{Strict lower bounds excluding $t\le.066$ when $\sqrt c<.44$.}
\label{elim:small-table}
\end{table}

\subsection{Bounding the remaining radii}

We next prove that
\begin{equation}\label{elim:tail-cap}
 r_j\le.42\qquad(j\ge4).
\end{equation}
Suppose instead that $r_4>.42$, and put
\[
 K=\frac{d\cdot.42}{s+.42}.
\]
For every $u\ge.42$ we have
\begin{equation}\label{elim:clipped-charge}
 \max\{0,e(u)\}\ge K(s^2-u^2).
\end{equation}
Indeed, when $u\le s$, use
\[
 e(u)=\frac{du}{s+u}(s^2-u^2)
 \quad\hbox{and}\quad \frac{u}{s+u}\ge\frac{.42}{s+.42}.
\]
When $u>s$, the right side of \eqref{elim:clipped-charge} is negative,
so nonnegativity proves the inequality. Applying this to the first four
radii and using the minimum radius on the others gives
\[
 \sum_{j=1}^4\delta_j
 \ge K\left(4s^2-\sum_{j=1}^4r_j^2\right)
 \ge K\bigl(4s^2-1+(m-4)t^2\bigr).
\]
For $j\ge5$, ordering and normalization give
$r_j\le1/\sqrt5$. By \eqref{elim:minimum-lower}, these radii lie
in $[.066,1/\sqrt5]$. Concavity of $e$ therefore gives
\[
 \delta_j\ge\min\{e(.066),e(1/\sqrt5)\}=e(.066).
\]
The last equality follows from Table~\ref{elim:four-table}, where
$e(.066)=.023996676$ for $A=.929$ and
$e(.066)=.024192696$ for $A=.884$.
Combining the charges and replacing $t$ by $.066$ proves
\begin{equation}\label{elim:four-charge}
 1-A\ge K\bigl(4s^2-1+(m-4)(.066)^2\bigr)+(m-4)e(.066).
\end{equation}
The right side increases with $m$, with increment
$K(.066)^2+e(.066)>0$. At the smallest permitted $m$, its excess
over $1-A$ is already positive by Table~\ref{elim:four-table}.
This contradiction proves \eqref{elim:tail-cap}.

\begin{table}[ht]
\centering
\begin{tabular}{c|c|c|c}
$A$ & smallest $m$ & \shortstack{Lower bound for\\ $e(1/\sqrt5)$} &
 \shortstack{Lower bound for the right side of\\ \eqref{elim:four-charge} minus $1-A$}\\
\hline
$.929$ & $5$ & $.02979364$ & $.00092980$\\
$.884$ & $6$ & $.03879364$ & $.02229774$
\end{tabular}
\caption{Strict lower bounds excluding four radii greater than $.42$.}
\label{elim:four-table}
\end{table}

Put $k=2$ when $m=5$ and $k=3$ when $m\ge6$. There are at least
$k$ indices after the first three, and their radii belong to $[t,.42]$.
Their charges are each at least $\min\{e(t),e(.42)\}$. Moreover,
\begin{equation}\label{elim:cap-budget}
 ke(.42)-(1-A)=
 \begin{cases}.0027688,&A=.929,\\.0184672,&A=.884.
 \end{cases}
\end{equation}
Hence \eqref{elim:budget} forces $ke(t)\le1-A$.
Let $t_0$ be the smaller root of $ke(u)=1-A$:
\begin{equation}\label{elim:tzero}
 t_0=\frac{.415-\sqrt{.415^2-4d(1-A)/k}}{2d}.
\end{equation}
Equation~\eqref{elim:cap-budget} places $.42$ strictly between the
two roots of this concave quadratic. Since $t\le.42$, the larger-root
alternative is excluded, and $t\le t_0$. The enclosures
\begin{equation}\label{elim:tzero-enclosures}
 \begin{aligned}
 .1070554<t_0<.1070555&\qquad(A=.929),\\
 .1176565<t_0<.1176566&\qquad(A=.884)
 \end{aligned}
\end{equation}
imply $t_0<.109$ in the first case and $t_0<.12$ in both cases.
Furthermore, $e(.42)>(1-A)/k\ge e(t)$, so each of the $k$ selected
indices contributes at least $e(t)$.

\subsection{The final contradiction}

We claim that
\begin{equation}\label{elim:third-versus-minimum}
 c>2.6t.
\end{equation}
Otherwise, \eqref{elim:third-radius} places the third radius in
$[.23,\sqrt{2.6t}]$. It is distinct from the $k$ indices just charged,
so concavity of $e$ and \eqref{elim:budget} imply
\begin{equation}\label{elim:last-budget}
 1-A\ge\min\{e(.23),e(\sqrt{2.6t})\}+ke(t).
\end{equation}
We check both possible endpoint charges on the entire interval
$[.066,t_0]$.

First, $e$ is increasing up to $t_0$, the smaller root in
\eqref{elim:tzero}. Thus
\[
 e(.23)+ke(t)-(1-A)
 \ge e(.23)+ke(.066)-(1-A)>0,
\]
by Table~\ref{elim:final-table}. Second, the function
\[
 F(t)=e(\sqrt{2.6t})+ke(t)
     =.415\sqrt{2.6t}-2.6dt+k(.415t-dt^2)
\]
is strictly concave for $t>0$, since
\[
 F''(t)=-\frac{.415\sqrt{2.6}}{4t^{3/2}}-2kd<0.
\]
Its minimum on $[.066,t_0]$ is therefore attained at an endpoint.
Table~\ref{elim:final-table} gives $F(.066)>1-A$ and $2.6t_0<s^2$.
The latter implies $e(\sqrt{2.6t_0})>0$; since $ke(t_0)=1-A$, it
also gives $F(t_0)>1-A$. These inequalities contradict
\eqref{elim:last-budget} and prove \eqref{elim:third-versus-minimum}.

\begin{table}[ht]
\centering
\begin{tabular}{c|c|c|c}
$A$ & \shortstack{Exact value of\\ $e(.23)+ke(.066)-(1-A)$} &
 \shortstack{Lower bound for\\ $F(.066)-(1-A)$} &
 \shortstack{Lower bound for\\ $s^2-2.6t_0$}\\
\hline
$.929$ & $.031234252$ & $.01522916$ & $.00546157$\\
$.884$ & $.013199488$ & $.00253590$ & $.01376450$
\end{tabular}
\caption{Positive endpoint margins for \eqref{elim:last-budget}.
The second column is exact; the last two columns give strict lower bounds.}
\label{elim:final-table}
\end{table}

Finally, \eqref{geo:second-constraint}, \eqref{elim:third-versus-minimum},
and $0<t<.12$ give
\[
 10(3/\pi)^2t^2
 \ge1.5(c-t^2)^2
 >1.5(2.6-.12)^2t^2.
\]
But
\[
 1.5(2.6-.12)^2-10(3/\pi)^2>.10667806>0.
\]
This contradiction rules out $C>B^2$ and completes the upper bound.
